% =============================================================
%  EXTRAS: cajas de guía, cajas de texto, código fuente,
%  diagramas de Gantt, tablas y enlaces a repositorios.
% =============================================================


% -- CAJAS DE GUÍA DE LA CÁTEDRA -------------------------------
% Se usan dentro de los capítulos para indicar qué espera la
% cátedra en cada sección. Desaparecen del PDF con la opción
% de clase  guias=false  (recomendado para la entrega final).
%
%   \begin{guia}  ... \end{guia}          qué evalúa la cátedra
%   \begin{ejemplo} ... \end{ejemplo}     ejemplo de cómo escribirlo
%   \begin{pendiente} ... \end{pendiente} recordatorio para el autor
%   \completar{texto}                     marcador en línea (siempre visible)

\RequirePackage[most]{tcolorbox}
\RequirePackage[fixed]{fontawesome5}

% paleta CASSIE para las cajas
\definecolor{guiaColor}{HTML}{1a7f37}        % verde (tip)
\definecolor{ejemploColor}{RGB}{90, 90, 90}  % gris
\definecolor{pendienteColor}{HTML}{000000}   % negro: marcadores \completar{} y caja pendiente
\definecolor{notaColor}{HTML}{0969da}        % azul (note)
\definecolor{advertenciaColor}{HTML}{9a6700} % ocre (warning)

\newtcolorbox{guia}[1][\USAL@guia]{
    enhanced, breakable,
    colback=guiaColor!5, colframe=guiaColor,
    fonttitle=\bfseries, title={\faInfoCircle\ #1},
    arc=1mm, boxrule=0.6pt, left=2mm, right=2mm, top=1mm, bottom=1mm,
    fontupper=\small,
}
\newtcolorbox{ejemplo}[1][\USAL@ejemplo]{
    enhanced, breakable,
    colback=ejemploColor!4, colframe=ejemploColor,
    fonttitle=\bfseries, title={\faLightbulb\ #1},
    arc=1mm, boxrule=0.6pt, left=2mm, right=2mm, top=1mm, bottom=1mm,
    fontupper=\small,
}
\newtcolorbox{pendiente}[1][\USAL@pendiente]{
    enhanced, breakable,
    colback=pendienteColor!6, colframe=pendienteColor,
    fonttitle=\bfseries, title={\faPencil*\ #1},
    arc=1mm, boxrule=0.6pt, left=2mm, right=2mm, top=1mm, bottom=1mm,
    fontupper=\small,
}
% (\texorpdfstring evita avisos de hyperref cuando se usa dentro de un título)
\newcommand{\completar}[1]{\texorpdfstring{\textcolor{pendienteColor}{[\,#1\,]}}{[#1]}}

% con guias=false, las tres cajas se redefinen para descartar su contenido
\ifusal@guias\else
    \RequirePackage{environ}
    \RenewEnviron{guia}[1][]{}
    \RenewEnviron{ejemplo}[1][]{}
    \RenewEnviron{pendiente}[1][]{}
\fi


% -- CAJAS DE TEXTO PARA EL CUERPO -----------------------------
% Estas SÍ quedan en el documento final. Uso:
%   \begin{nota}[Título opcional] ... \end{nota}
%   \begin{advertencia} ... \end{advertencia}
%   \begin{resumenbox}[Síntesis] ... \end{resumenbox}

\newtcolorbox{nota}[1][\USAL@nota]{
    enhanced, breakable,
    colback=notaColor!8, colframe=notaColor!8,
    borderline west={0.8mm}{0pt}{notaColor},
    sharp corners, boxrule=0pt, left=3mm, right=2mm,
    fonttitle=\bfseries, coltitle=notaColor, detach title,
    before upper={\tcbtitle\par\smallskip},
    title={#1},
}
\newtcolorbox{advertencia}[1][\USAL@advertencia]{
    enhanced, breakable,
    colback=advertenciaColor!8, colframe=advertenciaColor!8,
    borderline west={0.8mm}{0pt}{advertenciaColor},
    sharp corners, boxrule=0pt, left=3mm, right=2mm,
    fonttitle=\bfseries, coltitle=advertenciaColor, detach title,
    before upper={\tcbtitle\par\smallskip},
    title={\faExclamationTriangle\ #1},
}
\newtcolorbox{resumenbox}[1][\USAL@resumen]{
    enhanced, breakable,
    colback=mainColor!4, colframe=mainColor!4,
    borderline west={0.8mm}{0pt}{mainColor},
    sharp corners, boxrule=0pt, left=3mm, right=2mm,
    fonttitle=\bfseries, coltitle=mainColor, detach title,
    before upper={\tcbtitle\par\smallskip},
    title={#1},
}


% -- CÓDIGO FUENTE ---------------------------------------------
% Se usa el paquete listings (no necesita Python ni shell-escape).
%   \begin{lstlisting}[language=Python, caption={...}, label={lst:...}]
%   \lstinputlisting[language=SQL]{ruta/al/archivo.sql}
\RequirePackage{listings}
\RequirePackage{algorithm}
\RequirePackage{algpseudocode}

\definecolor{commentsColor}{RGB}{36, 161, 156}
\definecolor{numColor}{RGB}{71, 96, 114}
\definecolor{stringColor}{RGB}{78, 154, 6}
\definecolor{kwColor}{RGB}{207, 92, 0}
\definecolor{bgCodeColor}{RGB}{249, 249, 249}
\definecolor{codeGray}{HTML}{888A85}

\lstdefinestyle{codeSnippet} {
  backgroundcolor=\color{bgCodeColor},
  basicstyle=\ttfamily\small\color{black},
  keywordstyle=\color{kwColor},
  stringstyle=\color{stringColor},
  commentstyle=\color{codeGray}\itshape,
  numberstyle=\tiny\color{black},
  breakatwhitespace=false,
  breaklines=true,
  captionpos=b,
  keepspaces=true,
  numbers=left,
  numbersep=10pt,
  showspaces=false,
  showstringspaces=false,
  showtabs=false,
  frame=single,
  frameround=ttt,
  tabsize=2,
  xleftmargin=2em,
  xrightmargin=1em,
  extendedchars=true,
  inputencoding=utf8,
  literate={á}{{\'a}}1 {é}{{\'e}}1 {í}{{\'i}}1 {ó}{{\'o}}1 {ú}{{\'u}}1
           {ñ}{{\~n}}1 {Á}{{\'A}}1 {É}{{\'E}}1 {Í}{{\'I}}1 {Ó}{{\'O}}1
           {Ú}{{\'U}}1 {Ñ}{{\~N}}1 {¿}{{?`}}1 {¡}{{!`}}1,
}
\lstset{style=codeSnippet}
\lstdefinelanguage{JavaScript}{
  keywords={typeof, new, true, false, catch, function, return, null, catch, switch, var, let, const, if, in, while, do, else, case, break, async, await, import, export, from, class, extends, this},
  sensitive=true,
  comment=[l]{//}, morecomment=[s]{/*}{*/},
  morestring=[b]', morestring=[b]", morestring=[b]`,
}
\lstalias{TypeScript}{JavaScript}
\renewcommand{\lstlistingname}{\USAL@listing}
\renewcommand{\lstlistlistingname}{\USAL@listoflistings}


% -- DIAGRAMA DE GANTT (cronograma) ----------------------------
% Ver el ejemplo en capitulos/02-metodologia.tex
\RequirePackage{pgfgantt}
% nombres de mes en español para el calendario del Gantt
\ifx\usal@language\usal@spanishname
  \def\pgfcalendarmonthshortname#1{\ifcase#1\or ene\or feb\or mar\or abr\or may\or jun\or jul\or ago\or sep\or oct\or nov\or dic\fi}
  \def\pgfcalendarmonthname#1{\ifcase#1\or enero\or febrero\or marzo\or abril\or mayo\or junio\or julio\or agosto\or septiembre\or octubre\or noviembre\or diciembre\fi}
\fi


% -- TABLAS ----------------------------------------------------
% columnas de ancho automático para tabularx / xltabular:
%   L = izquierda, C = centrada, R = derecha
\newcolumntype{L}{>{\raggedright\arraybackslash}X}
\newcolumntype{C}{>{\centering\arraybackslash}X}
\newcolumntype{R}{>{\raggedleft\arraybackslash}X}
\renewcommand{\arraystretch}{1.15}

% celda de encabezado en negrita
\newcommand{\tb}[1]{\textbf{#1}}

% semáforo para matrices de riesgo (probabilidad x impacto)
\newcommand{\riesgoAlto}[1]{\cellcolor{red!25}#1}
\newcommand{\riesgoMedio}[1]{\cellcolor{orange!25}#1}
\newcommand{\riesgoBajo}[1]{\cellcolor{green!20}#1}


% -- FIGURAS DE EJEMPLO ----------------------------------------
% El paquete mwe provee example-image-a, -b, -c para maquetar
% antes de tener las figuras reales.
\RequirePackage{mwe}


% -- ENLACE A REPOSITORIO --------------------------------------
%   \prettygitlink{usuario/repositorio}
\newcommand{\prettygitlink}[1]{
    \begin{table}[H]
    \centering
    \renewcommand{\arraystretch}{1.2}
        \begin{tabular}{cl}
        \hline
        {\large\faGithub} & \href{https://github.com/#1}{#1} \\ \hline
        \end{tabular}
    \end{table}
}


% -- VARIOS ----------------------------------------------------
\RequirePackage{afterpage}
\RequirePackage[export]{adjustbox}
